% !Mode:: "TeX:UTF-8"

\chapter{实验设计与结果分析}

为了验证本文提出的基于依赖度的属性约简与规则提取算法的有效性，本章在威斯康星乳腺癌（WBC）数据集上进行实验。实验重点考察属性约简的效率、依赖度的变化，以及提取出的决策规则的医学可解释性。

\section{实验环境与数据集}

实验环境为 Python 3.x，核心算法基于 \texttt{pandas} 与 \texttt{numpy} 实现。
WBC数据集包含569个样本，每个样本有9个条件属性（如Clump\_thickness，Bare\_nuclei等），决策属性 \texttt{Class} 取值为2（良性）和4（恶性）。在剔除包含缺失值的样本后，将全部特征作为条件属性集 \(C\) 输入模型。

\section{属性约简结果分析}

利用第3章实现的基于依赖度的启发式约简算法，对包含9个条件属性的原始决策系统进行求解，结果如下：

\begin{enumerate}
	\item \textbf{核属性（Core）}：算法首先计算出 \texttt{Bare\_nuclei}（裸核）的依赖度不可替代，它是整个决策系统中唯一的核属性。
	\item \textbf{约简过程}：以 \texttt{Bare\_nuclei} 为起点，算法通过前向贪心搜索，依次加入了 \texttt{Uniformity\_of\_cell\_size}（依赖度提升至0.7862）、\texttt{Clump\_thickness}（依赖度提升至0.9854）和 \texttt{Bland\_chromatin}（依赖度提升至1.0000）。
	\item \textbf{约简结果}：最终得到的约简属性集包含 4 个特征：\texttt{['Bare\_nuclei', 'Uniformity\_of\_cell\_size', 'Clump\_thickness', 'Bland\_chromatin']}。
\end{enumerate}

\textbf{分析}：原始数据集的9个特征被成功约简至4个，且约简后系统的依赖度依然保持在 1.0000（与原始系统完全一致）。这说明约简掉的5个特征（如细胞形状、边缘粘附等）在当前的医学评分体系下，对于区分良恶性肿瘤是冗余的。该结果证明了算法在保持分类能力不变的前提下，有效降低了数据维度。

\section{决策规则提取与医学解释}

在约简后的决策表上，设定最小支持度阈值为 0.01，最小置信度阈值为 0.9，提取出的部分高置信度规则如表 \ref{tab:rules} 所示。

\begin{table}[htbp]
	\centering
	\caption{部分高置信度决策规则 (Confidence = 1.0)}
	\label{tab:rules}
	\begin{tabular}{llcc}
		\hline
		\textbf{规则编号} & \textbf{规则内容 (IF ... THEN ...)} & \textbf{支持度} & \textbf{覆盖样本数} \\
		\hline
		R1 & IF Bare\_nuclei=0 AND Uniformity=1 AND Clump=1 AND Bland=1 & 0.0615 & 42 \\
		R2 & IF Bare\_nuclei=0 AND Uniformity=1 AND Clump=1 AND Bland=2 & 0.0542 & 37 \\
		R3 & IF Bare\_nuclei=0 AND Uniformity=1 AND Clump=3 AND Bland=2 & 0.0512 & 35 \\
		\hline
	\end{tabular}
\end{table}

\textbf{医学可解释性分析}：
表 \ref{tab:rules} 中的规则后件均为 `Class = 2`（良性）。从规则前件可以清晰地看出，当患者的 \texttt{Bare\_nuclei}（裸核）评分为最低的 0，且 \texttt{Uniformity\_of\_cell\_size}（细胞大小均匀度）评分为 1 时，无论细胞厚度和染色质如何变化，肿瘤为良性的概率达到了惊人的 100\%。
这一规则提取结果与临床医学常识高度吻合：裸核消失和细胞大小均匀是正常细胞的特征，而恶性肿瘤细胞往往具有核异型性（裸核增多、大小不一）。粗糙集算法通过纯数据驱动的方式，自动发现了这一关键的医学诊断逻辑。

\section{本章小结}

本章在WBC数据集上验证了基于依赖度的属性约简与规则提取算法的实用性。实验表明，算法能够剔除冗余特征，将9维特征降至4维，并提取出具备 100\% 置信度且符合医学病理逻辑的 IF-THEN 规则。相比于传统黑盒机器学习模型（如随机森林或SVM），粗糙集方法不仅能够完成分类，还能提供清晰的决策依据，在医疗辅助诊断领域具有极高的应用价值。